\chapter{Background}
\label{chap:background}

% Reading guide: each section establishes one conceptual prerequisite and ends
% with an open question or unresolved tension. The gaps identified here are
% resolved in the Methods chapter.

\section{Imagery Rehearsal Therapy}
\label{sec:bg-irt}
% QUESTION THIS SECTION ANSWERS: What is IRT and why is rescripting the hard part?
%
% Three-stage protocol: recording, rescripting, rehearsal
%   Each stage has defined purpose and bounded appropriate therapist behaviours
% Rescripting: patient rewrites the nightmare to reduce distress and restore agency
%   This is the cognitive core -- the phase where the therapist must make active
%   strategic decisions (not just listen or summarise)
% Treatment fidelity instruments (ENACT, NIH BCC): define adherence gradations
%   absent / partial / implemented -- the clinical standard for measuring whether
%   a therapist delivers an intervention correctly
% Clinical evidence base for IRT effectiveness
%
% TENSION LEFT OPEN: Fidelity assessment exists for human therapists delivering
%   manualized protocols. No equivalent exists for LLM-based agents -- the
%   evaluation instruments assume a human clinician.
%
% [Example] Nightmare excerpt -> rescripted version with strategy labels overlaid
%   (cognitive_reframe, confrontation) -- what a valid rescripting move looks like

Nightmare disorder affects an estimated 2--5\% of the adult population
\parencite{YPHMANIU,HM68TTLK,Q24CN9LD} and shows strong comorbidity
with \gls{ptsd} and depression \parencite{34M7Y3F8}. For patients with recurrent nightmares, the distress
extends beyond the dream itself: sleep avoidance, daytime fatigue, and
heightened anxiety create a cycle that standard sleep hygiene cannot break.
\gls{irt} is the recommended treatment for this condition
\parencite{34M7Y3F8, FGA5H3ZD}.

\gls{irt} follows a structured three-stage protocol. In the \emph{recording}
stage, the patient describes the nightmare in detail. In the \emph{rescripting}
stage, the patient and therapist collaboratively transform the nightmare
narrative to reduce distress and increase mastery. In the
\emph{rehearsal} stage, the new narrative is consolidated and the patient
prepares to rehearse it before sleep.
\textcite{FGA5H3ZD} describe the rescripting instruction as deliberately
open-ended: patients are told to ``change the nightmare any way you wish.''

The rescripting stage is the cognitive core of the protocol. Unlike recording
(which requires listening) or rehearsal (which requires repetition), rescripting
demands that the therapist choose a specific rescripting strategy and guide
the patient through it. These choices determine the character of the
intervention. A therapist who helps the patient confront a nightmare threat
delivers a different treatment from one who introduces a protective figure
or reframes the threat as harmless.

Clinical research has identified distinct categories of rescripting
strategies. \textcite{25UR8BHJ} coded five types from combat veteran
rescripts: alternative endings (58\%), positive image insertion (23\%),
threat transformation (13\%), reality markers (10\%), and distancing (8\%).
\textcite{62ZMHBC3} proposed a complementary classification through the
Multidimensional Mastery Scale, which groups rescripting outcomes along
physical, social, environmental, emotional, mystical, and avoidance
dimensions. Both frameworks show that rescripting is not a single
technique but a family of distinct interventions. Which strategy the
therapist selects is therefore a meaningful clinical decision, and the
consistency of that decision across sessions is what this thesis measures.

Treatment fidelity instruments provide the clinical standard for measuring
whether a therapist delivers an intervention correctly. The \gls{irt}
literature draws on two established frameworks: the ENACT rating
scale \parencite{GYNNP2CC} and the NIH BCC recommendations
\parencite{ITT882KN}. Both grade therapist adherence on a ternary scale:
\emph{absent}, \emph{partial}, or \emph{implemented}. This gradation
captures clinically meaningful differences. A therapist who briefly mentions
a rescripting strategy without developing it (partial) delivers a different
intervention from one who fully guides the patient through it (implemented).

This ternary framework provides a scoring foundation that transfers directly
to computational evaluation (see \autoref{sec:meth-m3}). However, existing
fidelity instruments assume a human clinician. They have not been adapted for
\gls{llm}-based agents, where the therapist is a probabilistic text generator
rather than a trained professional.

% [Example] Nightmare excerpt -> rescripted version with strategy labels
%   overlaid (cognitive_reframe, confrontation)


\section{Treatment Fidelity in LLM-Based Interventions}
\label{sec:bg-llm-healthcare}
% QUESTION THIS SECTION ANSWERS: Why is therapeutic AI different from a chatbot?
%
% Scope: protocol adherence in manualized interventions, not the full
%   LLM-healthcare landscape. General clinical AI is referenced only where
%   it establishes failure modes relevant to stability measurement.
% The general vs. protocol distinction -- conceptual grounding:
%   Treatment fidelity research (Mowbray et al., TYX7W5ZM; NIH BCC, ITT882KN)
%   formalises what it means for a therapist to deliver an intervention correctly:
%   adherence (doing the prescribed things) and competence (doing them well)
%   Fidelity is only a meaningful concept when there is a protocol to adhere to --
%   it is undefined for a general assistant
%   For rule-based chatbots, adherence is built in by construction -- responses
%   are pre-authored and reviewed once
%   For LLM-based protocol agents, adherence becomes probabilistic -- it must
%   be measured across runs, not assumed from a single review
% Capability vs. reliability -- the core gap:
%   A model can produce high-quality individual responses while being clinically
%   unreliable across sessions -- capability benchmarks cannot detect this
% General LLM failure modes (hallucination, stage confusion, persona
%   inconsistency) exist but are outside this study's scope. This study
%   focuses on the probabilistic nature of LLM outputs creating drift
%   in identical therapeutic contexts.
% Regulatory framing: EU AI Act high-risk classification; Medical Device
%   Regulation auditability and documentation requirements
%   Data sovereignty as a legal constraint pushing toward EU-resident models
%
% TENSION LEFT OPEN: Stability is a precondition for clinical deployment, but
%   no existing evaluation framework measures it for protocol-driven LLM agents.
%   Prior therapeutic AI evaluation is largely qualitative, post-hoc, or
%   single-session -- it cannot detect cross-run inconsistency.
%
% [FIG] EU AI Act risk tier diagram -- where autonomous therapeutic AI sits

The deployment of \glspl{llm} in mental health has expanded from empathy
tools and symptom screeners to protocol-driven therapeutic agents.
\textcite{DTKNEPHQ} survey this trajectory and argue that \glspl{llm} could
fundamentally reshape behavioural healthcare delivery, from diagnostic support
to full psychotherapeutic interventions. Early systems such as \modelname{Woebot}
\parencite{YE8KQTNH} delivered \gls{cbt} through rule-based dialogue trees
where every response was pre-authored and reviewed before deployment. For
these systems, treatment fidelity is a design property: the chatbot cannot
deviate from its script.

\gls{llm}-based agents operate differently. Each response is a fresh sample
from the model's output distribution. A single model can produce
high-quality, protocol-adherent responses on one run and strategically
inconsistent ones on the next. This distinction is central to the treatment fidelity literature.
\textcite{TYX7W5ZM} define fidelity as the extent to which an
intervention follows its intended protocol, covering dimensions such as
adherence (doing the prescribed things) and quality of delivery (doing
them consistently). For rule-based chatbots, adherence is
built in by construction. For \gls{llm}-based protocol agents, adherence
becomes probabilistic and must be measured across runs rather than assumed
from a single review.

\gls{llm}-based agents also face general failure modes such as hallucination,
stage confusion, and persona inconsistency. These are well-documented
concerns for clinical AI but fall outside the scope of this study. The
present work focuses on a narrower problem: the probabilistic nature of
\gls{llm} outputs means that identical therapeutic contexts can produce
different outputs on each run. Three types of inconsistency can occur,
each with different clinical implications.

First, the model may select different rescripting strategies across runs.
As established in \autoref{sec:bg-irt}, strategy selection is a
meaningful clinical decision. A patient who receives confrontation-based
rescripting in one session and cognitive reframing in the next encounters
two different interventions. In a protocol where the therapist is expected
to follow consistent principles, this is a fidelity problem in the sense
of \textcite{TYX7W5ZM}: the model fails to adhere to a stable treatment
approach.

Second, even when the model selects the same strategy, the therapeutic
language may vary between runs. Some variation is expected and acceptable:
a competent human therapist naturally paraphrases while maintaining the
same strategic intent. But large variation in wording, even under the same
strategy, could indicate unstable delivery.

Mowbray's fidelity framework
motivates this concern: if consistent delivery matters for human
therapists, it matters at least as much for \gls{llm} agents whose
outputs vary with every run. This study does not measure fidelity to a
clinical standard directly. Instead, it measures \emph{consistency across
runs}: does the model plan the same techniques each time (Method~1), does
it word the response similarly (Method~2), and does it carry out what it
declared (Method~3).

Third, the model may declare one strategy but carry out another. If a
model states that it will use confrontation but actually delivers a
calming sensory exercise, there is a mismatch between intent and
execution. This internal incoherence is distinct from the first two types:
the model is neither selecting a wrong strategy nor varying its language,
but failing to do what it said it would do.

Each type of inconsistency requires a different measurement approach.
\autoref{chap:methods} develops one method for each.

The \gls{eu} AI Act and Medical Device Regulation both require
auditability for software used in clinical contexts
(\autoref{sec:intro-ai-therapy}), yet the type of stability described
above remains largely unmeasured. Prior
therapeutic AI evaluation is qualitative, post-hoc, or single-session.
\textcite{B97RI7GA} represent the closest prior work with their BOLT
framework, which assesses \gls{llm} therapist behaviour across 13
behavioural codes drawn from motivational interviewing. BOLT asks whether
the model does the right thing. The present study asks a different question:
whether the model does the \emph{same} thing across repeated runs. Existing
frameworks, including BOLT, evaluate single responses rather than the
distribution of responses across repeated trials.

% [FIG] EU AI Act risk tier diagram


\section{Stochastic Behaviour in Language Models}
\label{sec:bg-stochastic}
% QUESTION THIS SECTION ANSWERS: Where does instability come from mechanically?
%
% Temperature sampling: at T>0, token selection is probabilistic -- small
%   per-token variance compounds across a full therapeutic response
% Autoregressive generation: each token conditions on all previous tokens --
%   a single divergent token can cascade into a completely different strategy
% Clinical consequence: the same patient on two different days may receive
%   strategically inconsistent care -- this is not benign noise in a
%   manualized protocol
% Prior stability work in NLP:
%   Targets factual consistency, output diversity, or benchmark reproducibility
%   Not therapeutic strategy stability or protocol adherence
%
% TENSION LEFT OPEN: The mechanical source of instability is well understood,
%   but its impact on structured therapeutic decision-making has not been
%   quantified. How much does stochastic sampling cost in clinical consistency?
%
% [FIG] Sampling diagram -- same frozen history at T=0.0 vs. T=0.7 ->
%   diverging strategy distributions

At temperatures above zero, \gls{llm} token selection is probabilistic.
The softmax distribution over the vocabulary is divided by the temperature
parameter before sampling, so higher temperatures flatten the distribution
and lower temperatures sharpen it toward the most likely token. At $T = 0$,
the model selects the highest-probability token deterministically (greedy
decoding). At $T > 0$, tokens are sampled from the resulting distribution,
introducing variance into the output.

This per-token variance compounds through autoregressive generation. Each
token is conditioned on all preceding tokens, so a single divergent token
early in a response can cascade into a different therapeutic
strategy. A model that selects ``confront'' rather than ``imagine'' at one
token position may generate an entirely different rescripting approach in the
subsequent response. The clinical impact of stochastic sampling is therefore not proportional
to per-token entropy but can amplify through the response.

The clinical consequence follows directly: the same patient presenting the same
nightmare on two different days may receive strategically inconsistent care.
In a general-purpose assistant, this variability is benign or even
desirable. In a structured therapy protocol like \gls{irt}, where the
therapist is expected to apply consistent techniques across sessions, it
maps directly to the three types of inconsistency described in
\autoref{sec:bg-llm-healthcare}: variable strategy selection, variable
response wording, and plan-response mismatch.

Prior work on \gls{llm} stability addresses different questions than
therapeutic consistency. \textcite{BMA3X2SS} survey the consistency
landscape, covering whether models give the same answer to the same
question across different phrasings and contexts. \textcite{RPEDMF63}
measure the opposite problem: \glspl{llm} generating repetitive plot
structures in creative writing, producing too little variation rather
than too much. \textcite{54A9ZU2N} focus on benchmark reproducibility,
showing that evaluation scores fluctuate across repeated runs even with
fixed seeds and $T = 0.0$.

These lines of research share the concern for cross-run variation but
none of them ask whether a model makes the same \emph{clinical decisions}
when presented with identical patient contexts.

The mechanical source of instability is well understood. What remains
unquantified is its impact on structured therapeutic decision-making. How
much does temperature-driven variance cost in clinical consistency? This
question motivates the temperature sweep in the experimental design
(\autoref{sec:meth-conditions}), which tests five points from $T = 0.0$
(deterministic) to $T = 0.6$ (typical vendor default for open-weight
models).

% [FIG] Sampling diagram -- same frozen history at different temperatures


\section{Evaluation Metrics for Text Generation}
\label{sec:bg-metrics}
% QUESTION THIS SECTION ANSWERS: How do people measure text similarity, and
%   why is it hard for therapy?
%
% Surface metrics (BLEU, ROUGE): token overlap
%   Therapeutically equivalent paraphrases score near zero
%   Lexically similar but strategically divergent responses score high
%   Ruled out: therapeutic paraphrasing is the norm; surface metrics
%   penalise valid variation
% BERTScore (Zhang et al., 2020): contextual embedding similarity
%   Meaning equivalence independent of surface form
%   Multiple embedding model options with different strengths
%   Highest human correlation achieved with NLI-fine-tuned models
%   Open question: how well does this transfer to therapeutic language?
% LLM-as-judge paradigm:
%   Evaluates criteria no scalar metric can capture
%   (e.g. "does this response implement a confrontation strategy?")
%   Known biases: self-preference, position bias, verbosity bias
%   Mitigation strategies in the literature: cross-model judging,
%   deterministic decoding, structured rubrics
%
% TENSION LEFT OPEN: Embedding similarity captures semantic similarity but
%   cannot assess strategic intent. LLM judges can assess intent but introduce
%   their own reliability concerns. No single metric covers what clinical
%   stability requires -- a multi-level approach is needed, but none exists
%   for therapeutic AI.
%
% [FIG] ROUGE vs. BERTScore on a matched therapeutic sentence pair --
%   where surface overlap fails and why embedding similarity is needed

Measuring consistency between text outputs requires a similarity metric.
The choice of metric determines what kind of consistency is captured, and
standard options fail in different ways for therapeutic language.

Surface-level metrics such as BLEU and ROUGE work by counting how many
short word sequences (n-grams) appear in both texts, producing a score
between 0 and 1.

In therapeutic contexts, token overlap is misleading. A therapist
naturally varies their wording while keeping the same clinical intent, so
two responses that implement the same strategy through different phrasing
score low on BLEU/ROUGE despite being therapeutically equivalent.

BERTScore \parencite{PUYCKUYY} avoids this problem by comparing texts in
contextual embedding space rather than at the token level. It measures
whether two texts convey the same meaning regardless of surface wording.
The embedding model used in this study,
\modelname{DeBERTa-XLarge-MNLI} \parencite{8MQC2H2V}, achieves the
highest correlation with human similarity judgements ($r = 0.778$) among
available models. No validation specific to therapeutic language exists,
so this study uses its general-domain performance as the best available
approximation (\autoref{sec:meth-m2}).

BERTScore captures semantic similarity but cannot assess strategic intent.
A metric that reports high similarity between two responses does not reveal
whether both responses implement the same therapeutic strategy. The
\gls{llm}-as-judge paradigm addresses this gap. By instructing a separate
\gls{llm} to evaluate whether a response implements a declared strategy,
it becomes possible to assess criteria that no scalar metric can capture:
``does this response implement a confrontation strategy?'' requires
understanding both the strategy definition and the response content.

\gls{llm}-as-judge approaches carry known limitations. The judge's
scores are not ground truth but a computational estimate. The judge model
brings its own biases: self-preference (rating outputs from its own family
more favourably), position bias (higher scores for responses presented
first), and verbosity bias (rewarding longer responses regardless of
quality) \parencite{LVLA9LZU}. These can be mitigated through cross-model
judging, structured rubrics, and explicit justification requirements, but
they cannot be eliminated entirely. Any study using an \gls{llm} judge
must treat its scores as approximate and, where possible, validate them
against human ratings. The mitigation and validation strategy adopted in
this study is detailed in \autoref{sec:meth-m3}.

No single metric covers what therapeutic stability requires. Embedding
similarity captures semantic similarity but misses strategic intent.
\gls{llm} judges can assess intent but introduce their own reliability
concerns. Surface metrics capture neither. Existing therapeutic AI
evaluations rely on a single metric type, leaving the gap between semantic
equivalence and strategic intent unaddressed. This motivates the
three-method design presented in \autoref{chap:methods}.

% NOTE: Chain-of-Thought and Plan Faithfulness content folded into
% Section 3.4 (The Plan Mechanism) where the declared-intent resolution
% is introduced. Citations: Wei et al. 2022, Turpin et al. 2023,
% Lanham et al. 2023.
